\chapter{Methodology}


This chapter presents the complete methodology for implementing and evaluating the intelligent caching framework. Section~\ref{sec:theoretical} establishes the theoretical foundations, Section~\ref{sec:data} describes the dataset, Sections~\ref{sec:preprocessing}--\ref{sec:features} detail the signal processing and feature extraction pipeline, Section~\ref{sec:caching} explains the caching implementation, Section~\ref{sec:experiment_setup} outlines the experimental design, and Section~\ref{sec:evaluation_plan} defines the evaluation metrics.

\section{Theoretical Framework}
\label{sec:theoretical}

The proposed framework is grounded in three key concepts: configuration fingerprinting, a distinction between persistent storage and intelligent caching, and two orthogonal evaluation dimensions. This section presents both the currently implemented components and the planned extensions that form the core contribution of this thesis.

\subsection{Configuration Fingerprinting}

Configuration fingerprinting provides the foundation for cache key generation and reproducibility tracking. Each experiment configuration is uniquely identified by a cryptographic hash computed as:

\begin{equation}
\text{fingerprint} = \text{SHA256}(\text{canonical\_json}(\text{config}))[:32]
\end{equation}

The fingerprint uses 32 hexadecimal characters (128 bits), providing strong collision resistance suitable for large-scale experiment tracking. The configuration dictionary includes:

\begin{itemize}
    \item \textbf{Random seed:} Ensures deterministic behavior (e.g., \texttt{42})
    \item \textbf{Code version:} Git commit or semantic version (e.g., \texttt{v1.0}, \texttt{git:abc123})
    \item \textbf{Model configuration:} Algorithm name and hyperparameters (e.g., \texttt{xgboost}, \texttt{max\_depth: 6})
    \item \textbf{Feature configuration:} Base feature count, correlation threshold, and top-K selection
    \item \textbf{Held-out subject ID:} Critical for LOSO---prevents cross-fold contamination (e.g., \texttt{SC4001})
\end{itemize}

An example fingerprint: \texttt{a3f2b8c1d4e5f6a7b8c9d0e1f2a3b4c5}.

The fingerprint generation module is fully implemented, providing the \texttt{LOSOFingerprint} class with methods for generating fingerprints from configuration objects or dictionaries. The \texttt{generate\_cache\_key()} function provides a convenient interface that combines the fingerprint with the held-out subject ID (e.g., \texttt{a3f2b8c1...\_SC4001}) for direct use as cache file names.

\subsection{Configuration-Independent vs. Configuration-Dependent Caching}
\label{sec:storage_vs_caching}

A fundamental distinction in this framework is between \textit{configuration-independent caching} (artifacts that depend only on the dataset) and \textit{configuration-dependent caching} (artifacts that require fingerprint-based invalidation when experimental parameters change). While both avoid redundant computation, they serve different purposes:

\textbf{Configuration-Independent Cache} refers to artifacts that are computed once per dataset version and remain valid across all experimental variations. These artifacts do not depend on model hyperparameters, feature selection criteria, or cross-validation strategies. In this framework, the extracted feature matrices constitute configuration-independent caching---they are computed once per subject and reused across all experiments.

\textbf{Configuration-Dependent Cache} refers to artifacts whose validity depends on experimental configuration. These artifacts require fingerprint-based invalidation when parameters change. The LOSO model cache, which stores trained models keyed by configuration fingerprint and held-out subject, exemplifies configuration-dependent caching.

Table~\ref{tab:storage_hierarchy} summarizes this distinction:

\begin{center}
\colorbox{yellow!20}{%
\begin{minipage}{0.9\textwidth}
\textbf{Working Draft Note:} 
currently in last development Status and finding bugs + fixing table layout
\end{minipage}
}
\end{center}
\begin{table}[ht]
\centering
\begin{tabular}{|l|l|l|l|l|}
\hline
\textbf{Layer} & \textbf{Artifact} & \textbf{Size/Subject} & \textbf{Type} & \textbf{Status} \\
\hline
Layer 1 & 195 features (NPZ) & $\sim$1.1 MB & Persistent Storage & \textbf{Implemented} \\
\hline
Layer 2 & Trained LOSO models & $\sim$10--50 MB & Fingerprint Cache & \textbf{Implemented \& currently verfied} \\
\hline
\end{tabular}
\caption{Storage hierarchy distinguishing persistent storage from intelligent 
caching. Both layers are operational; verification testing demonstrated 
100\% cache hit rate and 4.5$\times$ speedup for identical experiment re-runs.}
\label{tab:storage_hierarchy}
\end{table}

\subsubsection{Layer 1: Feature Extraction as Persistent Storage (Implemented)}

The feature extraction layer computes \textbf{all 195 features} for each subject (representing the full 8-channel configuration: 6 EEG + EOG + EMG) and stores them in compressed NPZ format. This is classified as persistent storage because:

\begin{itemize}
    \item Features depend only on the raw EEG data and fixed preprocessing parameters
    \item Changing model hyperparameters or feature selection thresholds does not invalidate the stored features
    \item The full feature set is extracted once; downstream selection operates on this cached superset
\end{itemize}

\textbf{Feature Count Clarification.} The cache stores 195 features (23 per-channel features $\times$ 8 channels + 11 global features). For experiments using only the 6 EEG channels, features are filtered to 149 at query time ($<$1 second) via the \texttt{select\_channel\_features()} function, avoiding cache invalidation when switching between channel configurations.

This layer is fully implemented and provides measured speedups of approximately 224$\times$ for feature loading compared to re-extraction.

\subsubsection{Layer 2: LOSO Model Caching (Primary Contribution)}
\label{sec:loso_caching}

The LOSO model cache represents the \textbf{primary contribution} of this thesis. 
The complete system—including fingerprint generation, cache storage, retrieval logic, and automatic invalidation, is in implementation right now and validated through comprehensive testing (see Section~\ref{sec:verification}).



\begin{center}
\colorbox{yellow!20}{%
\begin{minipage}{0.9\textwidth}
\textbf{Working Draft Note:} 
As of today 13.01.2026 this is the current implementation status
\end{minipage}
}
\end{center}

\textbf{Implementation Status:} The caching system achieved 100\% cache hit rate 
in verification tests, with measured speedup of 4.5$\times$ for identical re-runs. 
Testing across 24 training configurations (3 subjects $\times$ 8 experimental 
combinations) demonstrated zero cache corruption errors and correct fingerprint-based 
invalidation when parameters changed.

\textbf{Why LOSO-Level Caching, Not Just Model Caching?}

A critical design decision is caching at the \textit{LOSO fold level} rather than caching a single global model. This distinction is essential:

\begin{enumerate}
    \item \textbf{Data Leakage Prevention:} Each LOSO fold has a unique fingerprint that includes the \texttt{held\_out\_subject} field. A model trained with Subject~A held out \textit{cannot} be retrieved when Subject~B is the test subject---the fingerprints differ (as demonstrated by the implemented \texttt{LOSOFingerprint} class), forcing a cache miss and retraining.
    
    \item \textbf{Scientific Validity:} In LOSO cross-validation, each fold produces a \textit{different} model because the training data differs. Caching a single ``global model'' would be scientifically invalid---there is no single model, but rather $N$ distinct models (one per subject held out).
    
    \item \textbf{Incremental Experimentation:} When adding new subjects or modifying experimental parameters, only affected folds require retraining. Cached models for unchanged configurations remain valid and are reused automatically.
\end{enumerate}

\textbf{Why Model-Level Caching, Not Result-Level Caching?}

Within each LOSO fold, the framework caches the \textit{trained model} rather than just the evaluation metrics:

\begin{itemize}
    \item \textbf{Post-hoc Analysis:} Cached models enable feature importance analysis, prediction on new data, or computation of additional metrics without retraining.
    \item \textbf{Reproducibility Verification:} Identical fingerprints guarantee identical model states, enabling verification that predictions (not just aggregate metrics) are reproducible.
    \item \textbf{Debugging Capability:} When unexpected results occur, the exact model state can be inspected rather than requiring reproduction of the entire training process.
\end{itemize}

The trade-off is increased storage (models are larger than metric summaries), but this is acceptable given the scientific benefits and modern storage capacities.

\subsection{Two Orthogonal Evaluation Dimensions}

The framework evaluation follows two orthogonal dimensions that structure the experimental analysis:

\begin{itemize}
    \item \textbf{WHY dimension (Optimization Goals):}
    \begin{itemize}
        \item Latency: Time to retrieve cached vs. computed results
        \item Throughput: Subjects processed per unit time
        \item Time Saved: Cumulative computation time avoided across experimental iterations
    \end{itemize}
    \item \textbf{WHERE dimension (Cached Artifact Levels):}
    \begin{itemize}
        \item Layer 1: Extracted feature matrices (persistent storage---implemented)
        \item Layer 2: Trained LOSO fold models (intelligent cache---in progress)
        \item Results: Final metrics (always recomputed from cached models for verification)
    \end{itemize}
\end{itemize}

This two-dimensional framework enables systematic evaluation of caching benefits across different experimental scenarios, from simple re-runs (maximum cache hits) to parameter sweeps (partial cache hits) to completely new configurations (cache misses establishing new baselines).

\section{Data: BOAS Dataset}

\label{sec:data}

\subsection{Dataset Overview}

This research uses the Bitbrain Open Access Sleep (BOAS) dataset, a publicly available polysomnographic EEG database. Table~\ref{tab:dataset_overview} summarizes the dataset characteristics.

\begin{table}[ht]
\centering
\begin{tabular}{|l|l|}
\hline
\textbf{Parameter} & \textbf{Value} \\
\hline
Dataset & Bitbrain Open Access Sleep (BOAS) \\
Total Subjects & 128 healthy adults \\
Epochs per Subject & $\sim$940 (varies by recording length) \\
Total Epochs & $\sim$120,000 \\
EEG Channels & 6 (F3, F4, C3, C4, O1, O2) following 10-20 system \\
Sampling Rate & 256 Hz (downsampled to 128 Hz after preprocessing) \\
Epoch Duration & 30 seconds \\
Samples per Epoch & 3,840 (at 128 Hz) \\
Sleep Stages & 5 (Wake, N1, N2, N3, REM) per AASM criteria \\
\hline
\end{tabular}
\caption{BOAS dataset characteristics.}
\label{tab:dataset_overview}
\end{table}

\subsection{Dataset Selection Justification}

The BOAS dataset was selected over Sleep-EDF for several reasons:
\begin{itemize}
\item \textbf{Public Availability:} Freely accessible via OpenNeuro, enabling reproducibility.
\item \textbf{Larger Subject Pool:} 128 subjects with consistent recording quality.
\item \textbf{Complete Channel Coverage:} All required EEG channels available across all subjects.
\item \textbf{Modern Annotations:} Sleep stages annotated according to AASM criteria.
\end{itemize}

\subsection{Expected Class Distribution}

Sleep stage classification exhibits inherent class imbalance, with N1 (light sleep) typically underrepresented. Based on literature and preliminary analysis, the expected distribution is approximately: Wake (15--20\%), N1 (5--10\%), N2 (45--50\%), N3 (15--20\%), REM (15--20\%). This imbalance will be addressed through macro-averaged metrics in evaluation.

\section{Data Validation and Integrity Verification}
\label{sec:data_validation}

Ensuring data quality and integrity is critical for reproducible sleep staging research. Prior to feature extraction, each polysomnographic recording undergoes a comprehensive 12-point validation protocol to detect common data issues that could compromise classification accuracy or introduce systematic biases.

\subsection{Motivation}

Sleep EEG datasets frequently contain artifacts from various sources: electrode disconnections during subject movement, signal clipping from amplifier saturation, flat channels from electrode detachment, or misalignment between signal data and annotation files. Without systematic validation, such issues may go undetected, leading to corrupted features that degrade classifier performance or, worse, produce misleadingly optimistic results when problematic epochs are inadvertently excluded from evaluation.

The BOAS dataset explicitly marks disconnection events (stage code 8) and artifacts (stage code -2) in the annotation files. However, additional integrity checks are necessary to verify temporal alignment between annotations and signal data, detect hardware failures not captured in annotations, and ensure data completeness.

\subsection{Validation Protocol}

The implemented validation protocol comprises 12 automated checks organized into four categories.

\textbf{Annotation Integrity (5 checks):}
\begin{itemize}
\item \textbf{Sufficient epochs:} $n_{\text{valid}} \geq 100$ ensures adequate data for statistical analysis
\item \textbf{All stages present:} $\{\text{W, N1, N2, N3, REM}\} \subseteq \text{recording}$ detects incomplete recordings or labeling errors
\item \textbf{Sufficient duration:} $t \geq 7200\text{s}$ (2 hours) excludes truncated or partial recordings
\item \textbf{Acceptable filter rate:} $\text{filtered}/\text{total} < 30\%$ indicates systematic problems when exceeded
\item \textbf{Temporal continuity:} $\max(\Delta t) \leq 1.5 \times \text{epoch\_duration}$ detects unexpected gaps in annotation timeline
\end{itemize}

\textbf{Structural Consistency (2 checks):}
\begin{itemize}
\item \textbf{Correct channels:} $n_{\text{channels}} \geq 6$ verifies required EEG montage is present
\item \textbf{Correct sampling rate:} $|f_s - f_{\text{target}}| < 1\text{ Hz}$ ensures consistent temporal resolution
\end{itemize}

\textbf{EDF-Annotation Alignment (2 checks):}
\begin{itemize}
\item \textbf{Annotations within EDF:} $\max(\text{onset}) + \text{epoch\_dur} \leq \text{EDF\_duration}$ prevents out-of-bounds data access
\item \textbf{Sample bounds valid:} $\max(\text{endsample}) \leq n_{\text{samples}}$ verifies sample indices reference existing data
\end{itemize}

\textbf{Signal Quality (3 checks):}
\begin{itemize}
\item \textbf{No invalid values:} $\nexists$ NaN $\vee$ $\pm\infty$ detects corrupted or missing samples
\item \textbf{No flat channels:} $\sigma_{\text{channel}} > 10^{-10}$ identifies dead electrodes or recording failures
\item \textbf{Reasonable amplitude:} $\max(|x|) < 10^6$ detects extreme values indicating unit errors or clipping
\end{itemize}

\subsection{Implementation}

Signal quality checks are performed on representative segments (first and last 30 seconds) rather than the complete recording to minimize computational overhead while still detecting common failure modes such as electrode detachment at recording start or end. For each validation failure, a descriptive error message is logged to facilitate debugging and enable transparent reporting of data exclusions.

\subsection{Validation Results: Preliminary Subset}

Table~\ref{tab:validation_results} shows validation results for the first three 
subjects, representing a pilot verification of the data quality protocol. 

\begin{table}[ht]
\centering
\begin{tabular}{|l|c|c|c|l|}
\hline
\textbf{Subject} & \textbf{Valid Epochs} & \textbf{Duration (h)} & \textbf{Filtered} & \textbf{Issues} \\
\hline
sub-1 & 914/915 & 7.63 & 0.1\% & 1 disconnection at $t=3690$s \\
\hline
sub-2 & 913/913 & 7.61 & 0.0\% & None \\
\hline
sub-3 & 997/997 & 8.31 & 0.0\% & None \\
\hline
\end{tabular}
\caption{Validation results for pilot dataset (3/128 subjects). Full dataset 
validation is ongoing; preliminary results confirm high data quality suitable 
for ML analysis.}
\end{table}

\textbf{Status:} As of this writing, comprehensive validation of all 128 subjects 
is in progress. The pilot results demonstrate that the BOAS dataset exhibits 
excellent signal integrity, with minimal artifact contamination ($<$0.1\% epoch 
exclusion rate) and high inter-rater agreement (83.5\% human-AI concordance). 
Complete validation results will be included in the final thesis submission.

All subjects passed the critical signal integrity checks (no NaN values, no flat channels, reasonable amplitudes), confirming that the BOAS dataset provides high-quality polysomnographic recordings suitable for machine learning analysis. The single disconnection event in subject 1 was correctly annotated in the original dataset and automatically excluded during preprocessing.


\subsection{Human vs.\ AI Label Agreement}

An additional quality metric available in the BOAS dataset is the agreement between human expert annotations (\verb|stage_hum|) and automated AI staging (\verb|stage_ai|). For subject 1, the inter-rater agreement was 83.5\%, which aligns with reported inter-scorer reliability in clinical polysomnography (typically 80--90\%). This provides a useful reference point for evaluating classifier performance: achieving accuracy significantly below human-AI agreement may indicate model issues, while exceeding it suggests the classifier captures patterns consistent with  the consensus human annotation standard.

\subsection{Importance for Reproducibility}

By systematically validating data integrity before analysis, this pipeline ensures that:

\begin{enumerate}
\item \textbf{No silent failures:} Data integrity issues are detected and 
      logged explicitly, preventing silent propagation of corrupted features through the pipeline
\item \textbf{Transparent exclusions:} Any filtered epochs or subjects are documented with specific reasons
\item \textbf{Consistent preprocessing:} All subjects undergo identical validation criteria
\item \textbf{Reproducible results:} The validation protocol can be applied to new datasets or by other researchers
\end{enumerate}

This approach aligns with best practices for transparent reporting in machine learning research and addresses concerns about hidden data preprocessing decisions that can affect reported classification accuracy.


\section{Preprocessing Pipeline}
\label{sec:preprocessing}

The preprocessing pipeline transforms raw EDF recordings into clean, epoched signals suitable for feature extraction.

\subsection{Signal Filtering}

The following signal conditioning is applied sequentially to each EEG channel:

\begin{enumerate}
\item \textbf{Bandpass Filter:} A 5th-order Butterworth filter with cutoff frequencies of 0.5--40 Hz removes DC drift and high-frequency noise while preserving sleep-relevant frequency bands.
\item \textbf{Notch Filter:} A notch filter at 50 Hz (Q=30) removes power line interference.
\item \textbf{Downsampling:} 256 Hz $\rightarrow$ 128 Hz using decimation with built-in anti-aliasing low-pass filter (prevents aliasing artifacts).
\end{enumerate}

All preprocessing parameters are included in the cache fingerprint, ensuring automatic invalidation when parameters change.

\subsection{Epoch Creation}

Continuous signals are segmented into 30-second epochs aligned with hypnogram annotations. Each epoch contains 3,840 samples (30 seconds $\times$ 128 Hz) across 6 channels, resulting in an array of shape $(n_\text{epochs}, 3840, 6)$ per subject.

\subsection{Quality Control}

Epochs are filtered according to the comprehensive validation criteria established in Section~\ref{sec:data_validation}. In summary, epochs must satisfy:
\begin{itemize}
\item Amplitude within $\pm$500 $\mu$V (reject artifacts)
\item No NaN or infinite values
\item Valid sleep stage annotation (exclude disconnection events, artifacts, and unknown stages)
\end{itemize}

The full 12-point validation protocol ensures data integrity before feature extraction begins.

\section{Feature Extraction: The 149-Feature Set}
\label{sec:features}

This research employs a comprehensive extraction of 149 EEG features per epoch strategy designed for robust cross-subject generalization. Features are computed for each 30-second epoch and organized into four domains, as detailed in Table~\ref{tab:feature_breakdown_149}.

\begin{table}[ht]
\centering
\begin{tabular}{|p{3.5cm}|p{1.5cm}|p{7cm}|}
\hline
\textbf{Domain} & \textbf{Count} & \textbf{Features} \\
\hline
Time-Domain & 60 & Mean, std, variance, min, max, peak-to-peak, RMS, skewness, kurtosis, zero-crossing rate (10 $\times$ 6 channels) \\
\hline
Frequency-Domain & 54 & Band powers (delta 0.5--4 Hz, theta 4--8 Hz, alpha 8--13 Hz, sigma 12--15 Hz, beta 13--30 Hz, gamma 30--40 Hz), spectral entropy, peak frequency, median frequency (9 $\times$ 6 channels) \\
\hline
Signal Complexity & 24 & Hjorth mobility, Hjorth complexity, Hurst exponent, Detrended Fluctuation Analysis (4 $\times$ 6 channels) \\
\hline
Inter-Channel & 11 & Coherence pairs (6), Phase Locking Value pairs (3), global entropy, global complexity \\
\hline
\textbf{Total} & \textbf{149} & \textbf{Comprehensive EEG-based feature set} \\
\hline
\end{tabular}
\caption{Complete 149-feature breakdown for sleep stage classification.}
\label{tab:feature_breakdown_149}
\end{table}

\subsection{Feature Engineering Strategy and Caching Prioritization}

The feature extraction pipeline implements a comprehensive 149-feature representation spanning time-domain statistics, frequency-domain band powers, complexity measures, and cross-channel connectivity metrics. This feature set was deliberately designed to balance discriminative power for sleep stage classification with computational tractability for caching optimization.

While additional features such as sample entropy, permutation entropy, band power ratios, or spindle density detection could marginally enhance classification performance, the current architecture intentionally prioritizes caching efficiency as the central research contribution. By establishing a fixed, comprehensive feature space, the system enables complete cache validity across experimental variations in feature selection, model hyperparameters, and cross-validation strategies---changes that would otherwise trigger costly recomputation.

This design philosophy reflects a strategic trade-off: rather than pursuing incremental gains in classification accuracy through feature proliferation, the framework optimizes for reproducibility, computational efficiency, and scalability. The result is a system where subsequent experiments achieve near-instantaneous feature retrieval, reducing pipeline latency from hours to seconds, minimizing computational resource consumption, and ensuring deterministic reproducibility across research iterations. Future extensions may incorporate additional feature domains; however, the current implementation demonstrates that intelligent caching of a well-designed feature set delivers substantial practical benefits that outweigh marginal accuracy improvements from feature expansion.

\subsection{Feature Selection Strategy}

After extracting all 149 features, optional feature selection is applied using a two-stage pipeline:

\textbf{Stage 1: Correlation-Based Filtering}
\begin{itemize}
\item Removes highly correlated features using Pearson correlation
\item Threshold options: 0.75, 0.90, or None (no filtering)
\item Prevents multicollinearity in downstream models
\end{itemize}

\textbf{Stage 2: ANOVA F-statistic Selection}
\begin{itemize}
\item Ranks features by class-discriminative power using F-statistic (\texttt{f\_classif})
\item Top-K options: 30, 50, or None (retain all 149 features)
\item ANOVA selected over Mutual Information due to approximately 76$\times$ faster computation with minimal accuracy difference (0.2\% in preliminary benchmarks: ANOVA 80.1\% vs.\ MI 80.3\%)\footnote{Timing benchmarks: ANOVA 0.078s vs.\ MI 5.956s for top-50 feature selection on the 5-subject pilot dataset (benchmark conducted December 2025).}
\end{itemize}

\textbf{Implementation Note:} Feature selection is performed \textbf{globally} (not per-fold) to enable consistent cache fingerprinting. This design choice prioritizes caching efficiency while accepting minimal impact on cross-validation purity. The Stage 2 cache stores all 149 features, and selection parameters only affect the final model training step.

\subsection{Feature Extraction Optimization}

To enable efficient processing of the BOAS dataset (128 subjects, $\sim$1000 epochs each, $\sim$128,000 total epochs), several optimization strategies were evaluated and implemented.

\subsubsection{Implemented Optimizations (No Quality Loss)}

\textbf{Multiprocess Parallelization:} Feature extraction was parallelized across CPU cores using Python's joblib library with process-based workers. Thread-based parallelization was initially attempted but provided no speedup due to Python's Global Interpreter Lock (GIL) blocking concurrent execution of CPU-bound NumPy/SciPy operations. Process-based parallelization achieved near-linear scaling, reducing per-epoch extraction time from $\sim$4.7s to $\sim$1.1s on an 8-core system (4.3$\times$ speedup).

\textbf{Optimized Complexity Measures:} The computationally expensive complexity features (Detrended Fluctuation Analysis, Hurst exponent, Hjorth parameters) were accelerated using the antropy library, which provides Numba JIT-compiled implementations of identical peer-reviewed algorithms. This optimization maintained mathematical equivalence while reducing complexity feature computation time by approximately 2--3$\times$.

\textbf{Feature Caching Strategy:} A comprehensive caching system was implemented to compute features once for all 6 EEG channels and store them in compressed NumPy archives. Subsequent experiments can load cached features and filter to the required channel subset at runtime, eliminating redundant computation when comparing channel configurations.

\subsubsection{Evaluated But Not Implemented (Would Reduce Quality)}

\textbf{Approximate Hurst Exponent:} Replacing the DFA-based Hurst exponent with faster auto\-correlation-based approximations was considered but rejected. While offering 5--10$\times$ speedup, this approach produces systematically different values (10--15\% deviation) and would compromise comparability with existing sleep staging literature that relies on standard DFA.

\textbf{Reduced DFA Scale Resolution:} Decreasing the number of window scales in DFA from 10 to 5 was evaluated. Although this would halve DFA computation time, it introduces 2--3\% variance in the resulting exponents, potentially affecting the discrimination of sleep stages with similar spectral profiles (particularly N2/N3 differentiation).

\textbf{Sample Entropy Approximation:} Fast approximate entropy measures were considered but not implemented, as sample entropy's $O(N^2)$ complexity was already addressed through parallelization, and approximate methods would alter the feature's discriminative properties for sleep stage classification.

\subsubsection{Performance Summary}

Table~\ref{tab:optimization_summary} summarizes the optimization decisions and their impact.

\begin{table}[ht]
\centering
\begin{tabular}{|l|c|c|c|}
\hline
\textbf{Optimization} & \textbf{Status} & \textbf{Speedup} & \textbf{Quality Impact} \\
\hline
Process parallelization & Implemented & 4.3$\times$ & None \\
\hline
Antropy complexity & Implemented & 2--3$\times$ & None \\
\hline
Feature caching & Implemented & N/A* & None \\
\hline
Approximate Hurst & Rejected & 5--10$\times$ & $-$10--15\% accuracy \\
\hline
Reduced DFA scales & Rejected & 1.5$\times$ & $-$2--3\% variance \\
\hline
\end{tabular}
\caption{Feature extraction optimization decisions. *Caching eliminates recomputation entirely for repeated experiments.}
\label{tab:optimization_summary}
\end{table}

The combined optimizations (process-based parallelization on 8 cores + Numba-accelerated complexity measures via antropy) reduced total feature extraction time from approximately 150 hours (sequential, unoptimized) to approximately 53 minutes for the complete 128-subject dataset. With the caching system enabled, subsequent runs retrieve cached features near-instantly ($<$1 minute), representing a 227$\times$ speedup over the cold start.



\section{Caching Framework Implementation}
\label{sec:caching}

\subsection{Fingerprint Generation}

Cache keys are generated using deterministic JSON serialization followed by SHA-256 hashing:

\begin{verbatim}
import hashlib
import json

def generate_fingerprint(config: dict) -> str:
    """Generate SHA-256 fingerprint from configuration."""
    canonical = json.dumps(config, sort_keys=True, default=str)
    return hashlib.sha256(canonical.encode()).hexdigest()[:32]
\end{verbatim}

\subsection{Cache Storage Structure}

The cache uses a directory structure organized by fingerprint:

\begin{verbatim}
cache/
+-- features/ # Extracted features (all 195, config-independent)
| +-- subject_sub-1_full.npz  # NO fingerprint - always valid
| +-- subject_sub-2_full.npz
| +-- ...
+-- loso_model_cache/ # Trained models (config-dependent)
  +-- {fingerprint}_sub-1.joblib  # WITH fingerprint
  +-- {fingerprint}_sub-2.joblib
  +-- cache_registry.json  # Fingerprint metadata
\end{verbatim}

\textbf{Critical Design Decision:} Layer 1 (features) uses filenames 
\textit{without} fingerprints because features are configuration-independent, they 
remain valid regardless of model hyperparameters or feature selection settings. 
Layer 2 (models) includes fingerprints in filenames because models become invalid 
when training configurations change. This naming convention prevents accidental 
cache misuse and enables visual inspection of cache validity.

\subsection{Cache Hit/Miss Logic}

The caching workflow follows this decision process:

\begin{enumerate}
\item Generate fingerprint from current configuration
\item Check if cached file exists with matching fingerprint
\item If \textbf{HIT}: Load cached data, record time saved
\item If \textbf{MISS}: Compute data, save to cache, record computation time
\end{enumerate}

\subsection{Current Implementation: Disk-Based Caching}


\begin{center}
\colorbox{yellow!20}{%
\begin{minipage}{0.9\textwidth}
\textbf{Working Draft Note:} 
rewriting this section right now still accurate here but need to adjust chapter 4 accordingly
\end{minipage}
}
\end{center}

The implemented caching system uses persistent SSD storage for all trained models. 
When a model is requested, the system checks for the existence of a cache file 
matching the configuration fingerprint and held-out subject identifier. If found, 
the model is loaded from disk; otherwise, it is trained and saved for future use.

\textbf{Lookup Strategy:}
\begin{enumerate}
\item Generate fingerprint from current configuration
\item Check if \verb|{fingerprint}_{subject}.joblib| exists on disk
\item If \textbf{HIT}: Load model from disk ($\sim$50ms)
\item If \textbf{MISS}: Train model ($\sim$300ms), save to disk, return
\end{enumerate}

This single-tier approach is well-suited to the experimental workflow of this 
thesis, where each model is typically accessed once per experimental run. The 
measured 4.5$\times$ speedup (Section~\ref{sec:verification}) demonstrates that 
disk-based caching alone provides substantial performance benefits for 
batch-oriented ML experimentation.





\begin{center}
\colorbox{yellow!20}{%
\begin{minipage}{0.9\textwidth}
\textbf{Working Draft Note:} 
delete here and put into future work?
\end{minipage}
}
\end{center}

\subsection{Future Extension: Two-Tier RAM+Disk Caching}

For production environments requiring interactive latency (e.g., real-time sleep 
analysis dashboards where the same subjects are analyzed repeatedly within minutes), 
the architecture could be extended with an in-memory cache layer:

\textbf{Two-Tier Lookup Strategy:}
\begin{enumerate}
\item \textbf{RAM Cache:} Check least-recently-used (LRU) cache in memory 
      ($\sim$0.1ms load time)
\item \textbf{Disk Cache:} If RAM miss, check persistent SSD storage ($\sim$50ms), 
      then store in RAM
\item \textbf{Compute:} If both miss, train model ($\sim$300ms), save to disk 
      and RAM
\end{enumerate}

\textbf{When Two-Tier Caching Provides Value:}
\begin{itemize}
\item \textbf{Interactive applications:} Dashboards where users repeatedly analyze 
      the same subjects with parameter adjustments
\item \textbf{High query frequency:} Systems processing $>$10 requests per minute 
      for the same models
\item \textbf{Low-latency requirements:} Applications requiring $<$100ms response times
\end{itemize}

\textbf{Why Not Implemented in This Thesis:}
The current experimental workflow accesses each cached model exactly once per run 
(128 LOSO folds $\times$ 1 access = 128 sequential lookups). Since there is no 
repeated access to the same model within a single experimental run, the RAM cache 
would never achieve hits beyond the initial disk load. Additionally, the 128 
cached models occupy $\sim$175 MB on disk but would require $\sim$700 MB of RAM 
if all were held in memory simultaneously, which could impact system stability 
during concurrent processing tasks.

For the batch-oriented experimentation conducted in this research, disk-based 
caching provides optimal cost-benefit trade-off: substantial speedup (4.5$\times$) 
with minimal memory overhead and complete persistence across sessions.

\textbf{Implementation Considerations for Future Work:}
A production implementation could use Python's \verb|functools.lru_cache| decorator 
or a dedicated caching library such as \verb|cachetools| to manage the RAM tier. 
The cache size would be tuned based on available memory (typically 10--50 models 
for a system with 4--8 GB available RAM). Complete implementation details and 
performance projections are documented in the project repository.

\subsection{Handling Data Changes and Cache Invalidation}

A critical challenge in ML experiment caching is ensuring cache validity when underlying data changes. Our implementation addresses this through a multi-layered integrity system embedded in the cache metadata. Each cached entry stores a configuration fingerprint---a SHA-256 hash of all preprocessing parameters, alongside explicit version identifiers (\texttt{data\_version} and \texttt{feature\_version}) and source file metadata.

When raw data modifications occur, three invalidation strategies are available:

\begin{enumerate}
\item \textbf{Version Bumping:} Incrementing the \texttt{data\_version} parameter in the configuration file causes the fingerprint to change, automatically invalidating all cached entries without requiring manual intervention. This approach preserves the cache structure while ensuring fresh computation.

\item \textbf{Cache Purging:} For significant data changes or version migrations, the global cache directory can be deleted entirely, triggering a complete recomputation on the next execution.

\item \textbf{Strict Validation:} An optional \texttt{strict\_validation} mode computes file checksums during cache lookup, detecting any byte-level changes to source files. While more computationally expensive, this provides cryptographic guarantees of data integrity.
\end{enumerate}

Our implementation defaults to version-based validation (Option 1), balancing integrity assurance with computational efficiency. The fingerprint system ensures that configuration drift---such as modified filter parameters or channel selections---automatically triggers recomputation, while unchanged configurations benefit from cached results. This approach aligns with reproducibility requirements in scientific computing, where identical inputs must yield identical outputs.



\section{Experimental Setup}
\label{sec:experiment_setup}

\subsection{Technical Environment}

Table~\ref{tab:technical_env} specifies the software environment for reproducibility. All code is developed and tested with Python 3.12.2 and the following dependencies on an HP ProBook x360 435 G7 laptop.

\begin{table}[ht]
\centering
\begin{tabular}{|l|l|l|}
\hline
\textbf{Component} & \textbf{Installed Version} & \textbf{Purpose} \\
\hline
Python & 3.12.2 & Core language \\
\hline
scikit-learn & 1.4.x+ & ML models (Random Forest, feature selection) \\
\hline
XGBoost & 3.0.2 & Gradient boosting classifier (primary model) \\
\hline
PyTorch & 2.7.0 & Deep learning framework for FNN model \\
\hline
MNE-Python & 1.9.0 & EEG data loading and signal processing \\
\hline
SciPy & 1.15.3 & Signal filtering, spectral analysis \\
\hline
NumPy & 2.1.3 & Numerical computations, array operations \\
\hline
Pandas & 2.2.3 & Data manipulation and organization \\
\hline
joblib & 1.5.1 & Parallel processing, model serialization \\
\hline
pickle & Built-in & Model/data serialization for caching \\
\hline
\textbf{Random Seed} & \textbf{42} & \textbf{All stochastic operations} \\
\hline
\end{tabular}
\caption{Technical environment and software versions for Sleep-EDF ML pipeline.}
\label{tab:technical_env}
\end{table}

\textbf{Hardware Environment:}

\begin{table}[ht]
\centering
\begin{tabular}{|l|l|}
\hline
\textbf{Component} & \textbf{Specification} \\
\hline
Machine & HP ProBook x360 435 G7 (Convertible Laptop) \\
\hline
Processor & AMD Ryzen 7 4700U with Radeon Graphics \\
 & (8 cores, 8 logical processors, 2000 MHz) \\
\hline
RAM (Installed) & 16.0 GB \\
\hline
RAM (Available during operation) & 1.6 GB \\
\hline
GPU & AMD Radeon Graphics (integrated) \\
 & No dedicated NVIDIA GPU \\
\hline
Storage & SSD on C: drive (primary boot device) \\
\hline
Operating System & Microsoft Windows 11 Pro (Build 26100) \\
\hline
System Type & x64-based PC with UEFI BIOS \\
\hline
\end{tabular}
\caption{Hardware environment specifications for Sleep-EDF ML pipeline development.}
\label{tab:hardware_env}
\end{table}

\textbf{Memory Management Strategy:}
With only 1.6 GB available RAM during operation, this pipeline employs aggressive memory optimization:
\begin{itemize}
	\item \textbf{Batch Processing:} Processes subjects in batches of 5--10 to keep feature matrices in memory
	\item \textbf{Memory Mapping:} Uses \texttt{numpy.memmap} for large arrays (loads only required portions into RAM on-demand)
	\item \textbf{Feature Caching:} Caches all 149 features to disk to avoid recomputation, reducing memory footprint of intermediate representations
	\item \textbf{Garbage Collection:} Explicit cleanup between subjects to free memory before next batch
	\item \textbf{Virtual Memory:} Leverages 15.4 GB page file if necessary, with acceptable I/O overhead
\end{itemize}

\textbf{Installation and Reproducibility:}
All dependencies are specified in the provided \texttt{requirements.txt} file. To replicate the development environment, run:

\begin{lstlisting}[language=bash]
pip install -r requirements.txt
\end{lstlisting}

\textbf{Performance Expectations:} This pipeline was developed and tested on a mobile workstation with 8 CPU cores and highly constrained available RAM (1.6 GB). Full LOSO cross-validation (128 folds) with feature caching requires approximately 12--16 hours of computation, with significant disk I/O overhead due to memory constraints. Results and timing metrics reported in Chapter~\ref{chap:results} are specific to this resource-constrained configuration. Users with more generous RAM availability will observe substantially faster execution, particularly for batch processing and model training phases.

\textbf{Note on Versions:} All installed versions are compatible and have been validated for this pipeline. For strict reproducibility in publication or deployment, consider pinning to exact versions in \texttt{requirements.txt}.
\subsection{Experimental Grid}

The experiments evaluate 18 configurations, as shown in Table~\ref{tab:exp_grid}.

\begin{table}[ht]
\centering
\begin{tabular}{|l|l|l|}
\hline
\textbf{Dimension} & \textbf{Options} & \textbf{Count} \\
\hline
Models & XGBoost, Random Forest & 2 \\
\hline
Correlation Thresholds & 0.75, 0.90, None & 3 \\
\hline
Top-K Features & 30, 50, None (all 149) & 3 \\
\hline
\textbf{Total Configurations} & 2 $\times$ 3 $\times$ 3 & \textbf{18} \\
\hline
\end{tabular}
\caption{Experimental grid: 18 configurations to evaluate.}
\label{tab:exp_grid}
\end{table}

Each configuration undergoes Leave-One-Subject-Out (LOSO) cross-validation with 128 folds, resulting in 2,304 total model training iterations.

\textbf{Note on FNN:} A Feedforward Neural Network (input$\rightarrow$64$\rightarrow$32$\rightarrow$5) 
is implemented and verified but excluded from final thesis experiments due to 
PyTorch dependency not being installed in the test environment. Preliminary 
results on 5 subjects achieved $\sim$78\% accuracy with successful caching 
(4.2$\times$ speedup on re-runs). The implementation is fully functional and 
could be included in future work with minimal additional effort ($<$1 hour 
computation time for full 18-configuration experiment once PyTorch is installed).




\begin{center}
\colorbox{yellow!20}{%
\begin{minipage}{0.9\textwidth}
\textbf{Working Draft Note:} 
as fnn is already partly implemented, but takes a lot of time to train (no dedicated gpu)fnn will not be used as of right now. maybe I am upgrading my pc soon to a laptop with dedicated gpu then i would like to incorporate the fnn model too
\end{minipage}
}
\end{center}


\subsection{Implementation Challenges and Resolutions}
\label{sec:implementation_challenges}


\begin{center}
\colorbox{yellow!20}{%
\begin{minipage}{0.9\textwidth}
\textbf{Working Draft Note:} 
note to myself, ask supervisor if these chapters should stay /or be canceld out - currently not sure about this chapters \end{minipage}
}
\end{center}

During development and testing, two critical bugs were identified and resolved. 
Documenting these issues provides insight into the subtleties of implementing 
ML experiment caching and may benefit researchers developing similar systems.

\subsubsection{Feature Mismatch Bug: Incomplete Fingerprinting}

\textbf{Symptom:} Models trained with 37 features crashed during loading when 
feature selection produced 33 features, with error: \texttt{X has 33 features, 
but XGBClassifier is expecting 37 features}.

\textbf{Root Cause:} The fingerprint generation included feature selection 
\textit{parameters} (correlation threshold, top-K count) but not the 
\textit{resulting feature names}. When different random seeds or data subsets 
produced different selected features under identical selection parameters, 
the fingerprint incorrectly indicated cache compatibility.

\subsubsection{Impact on Research Methodology}

These bugs were discovered through systematic testing before large-scale 
experiments, preventing potential corruption of results. The debugging process 
highlighted the importance of:
\begin{enumerate}
\item \textbf{Early verification testing} before committing to full experimental runs
\item \textbf{Explicit fingerprint design} that captures all sources of model variation
\item \textbf{Serialization-aware model design} for frameworks intended to support caching
\end{enumerate}

Both fixes are validated through the comprehensive test suite described in 
Section~\ref{sec:verification}.

\subsection{Cross-Validation Strategy}

\textbf{Leave-One-Subject-Out (LOSO)} cross-validation is the primary evaluation strategy:
\begin{itemize}
\item 128 folds (one per subject)
\item Each fold: 127 subjects for training, 1 subject for testing
\item Ensures true cross-subject generalization (no data leakage)
\item Total model trainings: 18 configurations $\times$ 128 folds = 2,304 trainings
\end{itemize}

\subsection{Caching Evaluation Scenarios}

Five scenarios will be executed to evaluate caching performance:

\begin{enumerate}
\item \textbf{First Run (Cold Start):} All cache misses, establishes baseline timing
\item \textbf{Identical Re-run:} All cache hits, measures maximum speedup
\item \textbf{Changed Preprocessing:} Invalidates Stage 1 and Stage 2 caches
\item \textbf{Changed Feature Selection:} Stage 1 and Stage 2 hit (selection at runtime)
\item \textbf{Add New Subjects:} Existing subjects hit, new subjects miss
\end{enumerate}

\section{Cache Verification Testing}
\label{sec:verification}

Prior to conducting the full 18-configuration experiment, the caching system 
underwent comprehensive verification to ensure correctness and measure performance.

\subsection{Test Design}


\begin{center}
\colorbox{yellow!20}{%
\begin{minipage}{0.9\textwidth}
\textbf{Working Draft Note:} 
maybe in final version cancel? as we use our real full grid + metrics and old test are not important anymore?
\end{minipage}
}
\end{center}

The verification protocol tested cache behavior across multiple scenarios using 
a reduced experimental grid (3 subjects $\times$ 8 configurations = 24 cache 
operations):

\begin{itemize}
\item \textbf{Subjects:} sub-1, sub-2, sub-3 (representative sample)
\item \textbf{Models:} XGBoost, Random Forest
\item \textbf{Feature Counts:} 30, 50, 149 (all)
\item \textbf{Correlation Thresholds:} 0.75, None
\item \textbf{Total Combinations:} $2 \times 3 \times 2 = 12$ configurations 
      (but $12 \times 2 = 24$ when counting both feature selection variants)
\end{itemize}

\textbf{Test Structure:}
\begin{enumerate}
\item \textbf{Cold Run:} Execute all 24 configurations with empty cache (expects 24 misses)
\item \textbf{Warm Run:} Re-execute identical configurations (expects 24 hits)
\item \textbf{Verification:} Compare predictions and metrics between runs
\end{enumerate}

\subsection{Verification Results}

Table~\ref{tab:cache_verification} shows the measured cache performance.

\begin{table}[ht]
\centering
\begin{tabular}{|l|c|c|c|c|c|}
\hline
\textbf{Run} & \textbf{Time} & \textbf{Hits} & \textbf{Misses} & \textbf{Hit Rate} & \textbf{Speedup} \\
\hline
Cold Start & 7.7s & 0 & 24 & 0\% & 1.0$\times$ (baseline) \\
\hline
Warm Re-run & 1.7s & 24 & 0 & 100\% & 4.5$\times$ \\
\hline
\end{tabular}
\caption{Cache verification results across 24 model training operations. Warm 
run achieved perfect cache hit rate with 4.5$\times$ speedup.}
\label{tab:cache_verification}
\end{table}

\subsection{Correctness Validation}

Beyond performance, the test suite verified cache correctness:

\begin{itemize}
\item \textbf{Prediction Identity:} Cached models produced bitwise-identical  predictions to freshly trained models (verified via \verb|numpy.array_equal()|)
\item \textbf{Metric Reproducibility:} Accuracy, F1, and kappa scores matched 
      to floating-point precision ($<10^{-10}$ difference)
\item \textbf{Zero Corruption:} All 24 pickle save/load cycles completed without errors
\item \textbf{Fingerprint Uniqueness:} Different configurations generated distinct 
      fingerprints (no hash collisions)
\end{itemize}

\subsection{Invalidation Testing}

Cache invalidation was verified by modifying configurations:

\begin{itemize}
\item \textbf{Changed Hyperparameter:} Modifying \verb|max_depth| from 6 to 8 
      $\rightarrow$ fingerprint changed $\rightarrow$ cache miss (correct)
\item \textbf{Changed Feature Selection:} Changing top-K from 30 to 50 $\rightarrow$ 
      fingerprint changed $\rightarrow$ cache miss (correct)
\item \textbf{Changed Held-Out Subject:} Testing sub-1 vs.\ sub-2 with identical 
      parameters $\rightarrow$ different fingerprints $\rightarrow$ no cross-fold 
      contamination (correct)
\end{itemize}

\subsection{Implications for Full Experiment}

These verification results establish confidence that:
\begin{enumerate}
\item The caching system correctly identifies cache hits/misses
\item Cached models are functionally identical to freshly trained models
\item Fingerprinting prevents inappropriate cache reuse across configurations
\item The measured 4.5$\times$ speedup on 24 operations extrapolates to substantial 
      time savings across the full 2,304-operation experiment (18 configs $\times$ 
      128 LOSO folds)
\end{enumerate}

The full experimental results (Chapter 5) will report cache performance across 
the complete dataset and experimental grid.


\section{Evaluation Plan}



\begin{center}
\colorbox{yellow!20}{%
\begin{minipage}{0.9\textwidth}
\textbf{Working Draft Note:} 
eventuell metriken leicht anpassen
\end{minipage}
}
\end{center}
\label{sec:evaluation_plan}

\subsection{Primary Metrics: Caching Efficiency}

The primary thesis contribution is evaluated through caching metrics:

\begin{itemize}
\item \textbf{Cache Hit Rate:} $\frac{\text{hits}}{\text{hits} + \text{misses}}$ per stage
\item \textbf{Time Saved:} $\sum(\verb|expected_compute_time| - \text{cache\_load\_time})$
\item \textbf{Speedup Factor:} $\frac{\text{cold\_start\_time}}{\text{cached\_run\_time}}$
\item \textbf{Storage Overhead:} Total cache size in MB/GB
\end{itemize}

\subsection{Secondary Metrics: Classification Performance}

Classification metrics validate that the caching framework produces correct results:

\begin{itemize}
\item \textbf{Accuracy:} Overall correct classifications
\item \textbf{Macro F1-Score:} Mean F1 across all classes (handles imbalance)
\item \textbf{Cohen's Kappa:} Agreement beyond chance
\item \textbf{Per-Class F1:} Identifies problematic classes (especially N1)
\item \textbf{Confusion Matrix:} 5$\times$5 aggregated across all folds
\end{itemize}

\subsection{Reproducibility Validation}

Reproducibility will be verified through:

\begin{itemize}
\item \textbf{Fingerprint Determinism:} Same config $\rightarrow$ same fingerprint (verified across multiple runs)
\item \textbf{Result Consistency:} Same fingerprint $\rightarrow$ identical metrics (within floating-point tolerance)
\item \textbf{Resumability Test:} Interrupt pipeline $\rightarrow$ resume from cache $\rightarrow$ identical final results
\end{itemize}

\chapter{Expected Results}
